% Бүлэг 3

\chapter{Системийн шаардлага}
\label{Chapter3}
\pagecolor{white}

\section{Системийн ерөнхий тодорхойлолт}

ViralAI систем нь хэрэглэгч зөвхөн сэдэв оруулахад скрипт үүсгэх, зураг үүсгэх, дуу хоолой синтез хийх, видео угсрах гэсэн 4 үндсэн үе шатыг автоматаар гүйцэтгэж богино хэмжээний видео контент бэлдэдэг цогц веб систем юм. Система нь гурван үндсэн актортой: Хэрэглэгч (User), Контент Менежер (Content Manager), Супер Админ (Super Admin).

\section{Хэрэглэгчийн шаардлага}

\begin{table}[h!]
\centering
\caption{Хэрэглэгчийн шаардлагын жагсаалт}
\label{tab:user_req}
\begin{tabular}{|c|p{4cm}|p{7cm}|}
\hline
\textbf{ID} & \textbf{Шаардлага} & \textbf{Тайлбар} \\
\hline
UR-01 & Бүртгэл, нэвтрэлт & Хэрэглэгч бүртгүүлж, нэвтрэх боломжтой байна \\
\hline
UR-02 & Видео үүсгэх & Сэдэв оруулахад автоматаар видео үүсгэгдэнэ \\
\hline
UR-03 & Явц хянах & Видео үүсгэлтийн явцыг бодит цагт харна \\
\hline
UR-04 & Studio засварлагч & Видеогоо scene бүрт засварлах боломжтой \\
\hline
UR-05 & Видео жагсаалт & Өөрийн бүх видеог харах, удирдах \\
\hline
UR-06 & Projects & Видеог folder бүлгээр зохион байгуулах \\
\hline
UR-07 & Subscription & Сарын хязгаар болон Pro эрхийг удирдах \\
\hline
UR-08 & Admin панел & Бүх хэрэглэгч, видеог харж удирдах (Admin) \\
\hline
UR-09 & Статистик & Систем, хэрэглэгчийн ерөнхий статистик (Admin) \\
\hline
UR-10 & Аюулгүй байдал & JWT-д суурилсан баталгаажуулалт \\
\hline
\end{tabular}
\end{table}

\section{Функциональ шаардлага}

\begin{table}[h!]
\centering
\caption{Функциональ шаардлагын жагсаалт}
\label{tab:func_req}
\begin{tabular}{|c|p{3cm}|p{7.5cm}|}
\hline
\textbf{ID} & \textbf{Нэр} & \textbf{Тайлбар} \\
\hline
FR-01 & Бүртгэл & Хэрэглэгч email, нууц үгээр бүртгүүлнэ \\
\hline
FR-02 & Нэвтрэлт & JWT access + refresh token олгоно \\
\hline
FR-03 & Скрипт үүсгэлт & Groq/Claude API-аар scene бүрийн narration үүсгэнэ \\
\hline
FR-04 & Зураг үүсгэлт & Gemini Flash Image 2.5-аар scene зураг үүсгэнэ \\
\hline
FR-05 & Аудио үүсгэлт & ElevenLabs/Chimege/Gemini TTS-аар дуу хоолой синтез хийнэ \\
\hline
FR-06 & Видео угсралт & FFmpeg-ээр зураг+аудио+subtitle нэгтгэж MP4 гаргана \\
\hline
FR-07 & Явц дамжуулалт & SSE-ээр бодит цагт явцыг хэрэглэгчид мэдэгдэнэ \\
\hline
FR-08 & Cloudinary хадгалалт & Бэлэн видеог CDN-д байршуулна \\
\hline
FR-09 & Studio засварлагч & Scene тус бүрийн зураг/аудио/текст дахин үүсгэнэ \\
\hline
FR-10 & Re-render & Засварласан scene-г дахин угсарна \\
\hline
FR-11 & RBAC & Free/Pro/Admin эрхийн ялгаатай хяналт \\
\hline
FR-12 & Projects & Folder үүсгэх, видеог зөөх, устгах \\
\hline
FR-13 & Admin удирдлага & Хэрэглэгч, видео харах, устгах, статистик \\
\hline
FR-14 & Rate limiting & API хэтрүүлэлтийг хязгаарлана \\
\hline
\end{tabular}
\end{table}

\section{Функциональ бус шаардлага}

\begin{table}[h!]
\centering
\caption{Функциональ бус шаардлагын жагсаалт}
\label{tab:nonfunc_req}
\begin{tabular}{|c|p{3cm}|p{7.5cm}|}
\hline
\textbf{ID} & \textbf{Ангилал} & \textbf{Шаардлага} \\
\hline
NFR-01 & Гүйцэтгэл & 60 секундын видеог 5 минутын дотор үүсгэнэ \\
\hline
NFR-02 & Хандалт & API response P95 < 500ms байна \\
\hline
NFR-03 & Аюулгүй байдал & JWT, bcrypt, HTTPS, rate limiting ашиглана \\
\hline
NFR-04 & Найдвартай байдал & 99.5\%+ uptime, fallback AI strategy \\
\hline
NFR-05 & Өргөтгөх чадвар & Горизонталь scale хийх боломжтой архитектур \\
\hline
NFR-06 & Засварлах чадвар & Layered architecture, TypeScript, unit test \\
\hline
NFR-07 & Баримтжуулалт & Swagger API doc, TypeDoc, хэрэглэгчийн гарын авлага \\
\hline
NFR-08 & Хэрэглэгчийн туршлага & SPA, responsive дизайн, бодит цагийн явц \\
\hline
\end{tabular}
\end{table}

\section{Системийн хэрэглэгчийн төрлүүд}

Систем нь гурван үндсэн хэрэглэгчийн төрлийг дэмжинэ:

\begin{enumerate}
\item \textbf{Хэрэглэгч (User)} --- бүртгэлтэй хэрэглэгч. Free болон Pro гэсэн хоёр дэд ангилалтай. Free хэрэглэгч сард хязгаарлагдсан тооны видео үүсгэж зурагны re-render хийх боломжгүй. Pro хэрэглэгч илүү олон видео үүсгэж, бүх Studio боломжийг ашиглана.

\item \textbf{Контент Менежер (Content Manager)} --- системийн контентыг удирдах эрхтэй дунд шатны эрхлэгч.

\item \textbf{Супер Админ (Super Admin)} --- бүх хэрэглэгч, видео, статистикийг харж удирдах хамгийн өндөр эрхтэй хэрэглэгч.
\end{enumerate}

\begin{table}[h!]
\centering
\caption{Хэрэглэгчийн эрхийн матриц}
\label{tab:rbac}
\begin{tabular}{|l|c|c|c|c|}
\hline
\textbf{Үйлдэл} & \textbf{Free} & \textbf{Pro} & \textbf{Manager} & \textbf{Admin} \\
\hline
Видео үүсгэх & \checkmark (хязгар) & \checkmark & \checkmark & \checkmark \\
\hline
Studio засварлагч & \checkmark (хэсэгчилсэн) & \checkmark & \checkmark & \checkmark \\
\hline
Зургийн re-render & \texttimes & \checkmark & \checkmark & \checkmark \\
\hline
Projects удирдах & \checkmark & \checkmark & \checkmark & \checkmark \\
\hline
Бүх хэрэглэгч харах & \texttimes & \texttimes & \texttimes & \checkmark \\
\hline
Бүх видео харах & \texttimes & \texttimes & \texttimes & \checkmark \\
\hline
Статистик харах & \texttimes & \texttimes & \texttimes & \checkmark \\
\hline
\end{tabular}
\end{table}

\section{Use Case-уудын жагсаалт}

Системийн үндсэн Use Case-ууд:

\begin{itemize}
\item UC-01: Видео үүсгэх
\item UC-02: Бүртгэл, нэвтрэлт
\item UC-03: Видео үүсгэлтийн явц хянах (SSE)
\item UC-04: Скрипт үүсгэлт
\item UC-05: Видео болон медиа хадгалалт
\item UC-06: Subscription удирдах
\item UC-07: Studio засварлагчаар scene засварлах
\item UC-08: Хэсэгчилсэн re-render хийх
\item UC-09: Projects зохион байгуулах
\item UC-10: Admin удирдлага
\end{itemize}

\section{Use Case диаграм}

Дараах диаграмд системийн гурван үндсэн актор --- Хэрэглэгч, Контент Менежер, Супер Админ --- нарын хийх боломжтой бүх үйлдлүүдийг мөн гадаад AI үйлчилгээнүүдтэй (Groq/Gemini, ElevenLabs) харилцааг дүрслэв.

\begin{figure}[h!]
\centering
% \includegraphics[width=0.9\textwidth]{Figures/usecase_diagram.png}
\caption{ViralAI системийн Use Case диаграм}
\label{fig:usecase}
\end{figure}

\section{Use Case тодорхойлолтууд}

\subsection{Үндсэн үйл ажиллагааны Use Case-ууд}

\begin{table}[h!]
\centering
\caption{UC-01 --- Видео үүсгэх}
\label{tab:uc01}
\begin{tabular}{|l|p{9cm}|}
\hline
\textbf{Use Case ID} & UC-01 \\
\hline
\textbf{Нэр} & Видео үүсгэх \\
\hline
\textbf{Актор} & Хэрэглэгч (Free/Pro) \\
\hline
\textbf{Зорилго} & Хэрэглэгч сэдэв оруулахад систем автоматаар скрипт, зураг, дуу, видео бүтээнэ \\
\hline
\textbf{Trigger} & Хэрэглэгч ``Видео үүсгэх'' товчийг дарна \\
\hline
\textbf{Өмнөх нөхцөл} & Хэрэглэгч системд бүртгэлтэй бөгөөд нэвтэрсэн байна \\
\hline
\textbf{Үндсэн урсгал} &
1. Хэрэглэгч сэдэв, жанр, урт, хэл, зурагны стиль оруулна \newline
2. Систем Groq API-аар скрипт үүсгэнэ (UC-04) \newline
3. Систем Gemini-ээр scene бүрт зураг үүсгэнэ \newline
4. Систем ElevenLabs-аар дуу хоолой синтез хийнэ \newline
5. Систем FFmpeg-ээр зураг, дуу, текстийг нэгтгэн MP4 үүсгэнэ \newline
6. Систем видеог Cloudinary-д байршуулна (UC-05) \newline
7. Хэрэглэгчид холбоос болон мэдэгдэл ирнэ \newline
8. Хэрэглэгч Studio засварлагч руу шилжих сонголттой болно \\
\hline
\textbf{Альтернатив урсгал} &
2a. Groq timeout болвол Gemini API-д fallback хийнэ \newline
5a. FFmpeg алдаа гарвал retry логик ажиллана \\
\hline
\textbf{Дараах нөхцөл} & Видео амжилттай үүсгэгдэж, хэрэглэгчийн жагсаалтад харагдана \\
\hline
\end{tabular}
\end{table}

\begin{table}[h!]
\centering
\caption{UC-07 --- Studio засварлагчаар scene засварлах}
\label{tab:uc07}
\begin{tabular}{|l|p{9cm}|}
\hline
\textbf{Use Case ID} & UC-07 \\
\hline
\textbf{Нэр} & Studio засварлагчаар scene засварлах \\
\hline
\textbf{Актор} & Хэрэглэгч (Pro), Контент Менежер \\
\hline
\textbf{Зорилго} & Видеоны тодорхой scene-ийн narration, зураг, дуу хоолойг засварлана \\
\hline
\textbf{Үндсэн урсгал} &
1. Хэрэглэгч Studio засварлагч нээнэ \newline
2. Scene timeline-аас засварлах scene-г сонгоно \newline
3. Narration текстийг засварлах, эсвэл AI-аар дахин бичүүлнэ \newline
4. Зурагны prompt засварлаж дахин үүсгүүлнэ \newline
5. Дуу хоолойг дахин синтез хийнэ \newline
6. Preview-г харна \newline
7. Reassemble дарж эцсийн видео дахин угсарна \\
\hline
\textbf{Дараах нөхцөл} & Засварласан scene бүхий шинэ видео үүснэ \\
\hline
\end{tabular}
\end{table}

\section{Бүлгийн дүгнэлт}

Энэ бүлэгт системийн гурван үүрэгтэй хэрэглэгч --- Хэрэглэгч, Контент Менежер, Супер Админ --- нарын 10 хэрэглэгчийн шаардлага (UR), 14 функциональ шаардлага (FR), 8 функциональ бус шаардлага (NFR) болон 10 Use Case-г тодорхойлсон. RBAC эрхийн матрицыг Free/Pro/Admin гэсэн ялгаатайгаар тодорхойлж, Studio засварлагчийн Use Case урсгалыг дэлгэрэнгүүй гаргав.
